\documentclass[11pt]{article}

\usepackage[a4paper, margin=1in]{geometry}

\usepackage{amsmath}
\usepackage{amsfonts}
\usepackage{amssymb}
\usepackage{mathtools}
\usepackage{microtype}
\usepackage{lmodern}
\usepackage{graphicx}
\usepackage{float}
\usepackage{subcaption}
\usepackage{xcolor}
\usepackage{hyperref}

\hypersetup{
colorlinks=false,
linkcolor=blue,
urlcolor=blue,
citecolor=blue
}

\usepackage{fancyhdr}
\pagestyle{fancy}
\fancyhf{}
\lhead{3D Graphics Pipeline}
\rhead{\thepage}
\renewcommand{\headrulewidth}{0.4pt}

\usepackage{titlesec}

\titleformat{\section}
{\large\bfseries}
{\thesection}{1em}{}

\titleformat{\subsection}
{\normalsize\bfseries}
{\thesubsection}{1em}{}

\usepackage{listings}
\lstset{
basicstyle=\ttfamily\small,
breaklines=true,
frame=single
}

\newcommand{\R}{\mathbb{R}}
\newcommand{\vct}[1]{\mathbf{#1}}
\newcommand{\mat}[1]{\mathbf{#1}}

\title{\textbf{From Bits to Pixels} \\
\large The Mathematical and Physical Foundations of elementary 3D Graphics}
\author{Jack Wetherell}
\date{}

% ================================
% Document Start
% ================================
\begin{document}

\maketitle
\tableofcontents
\newpage

\section{Foundations: Numbers, Memory, and Representation}

\subsection{Bits, Bytes, and Binary Storage}

At the lowest level, all data in a computer is represented using binary digits.

\begin{itemize}
    \item A \textbf{bit} is a single binary value:
    \begin{equation}
        b \in \{0, 1\}
    \end{equation}

    \item A \textbf{byte} consists of 8 bits:
    \begin{equation}
        1\ \text{byte} = 8\ \text{bits}
    \end{equation}
\end{itemize}

A sequence of bits can represent integers using positional notation in base 2:

\begin{equation}
    N = \sum_{i=0}^{n-1} b_i \cdot 2^i
\end{equation}

where:
\begin{itemize}
    \item $b_i$ is the $i$-th bit (0 or 1)
    \item $i = 0$ corresponds to the least significant bit (LSB)
\end{itemize}

\paragraph{Example:}
\begin{equation}
    1011_2 = 1 \cdot 2^3 + 0 \cdot 2^2 + 1 \cdot 2^1 + 1 \cdot 2^0 = 11_{10}
\end{equation}

Memory is therefore a linear array of bytes:

\begin{equation}
    \text{Memory} = [\text{byte}_0, \text{byte}_1, \dots, \text{byte}_N]
\end{equation}

Each byte has an address:

\begin{equation}
    \text{Address} \in \mathbb{N}
\end{equation}


\subsection{Floating Point Representation (IEEE 754)}

Real numbers are represented approximately using the IEEE 754 floating point format.

A 32-bit floating point number (single precision) is composed of:

\begin{itemize}
    \item 1 bit for the sign ($s$)
    \item 8 bits for the exponent ($e$)
    \item 23 bits for the mantissa (fraction) ($m$)
\end{itemize}

The value represented is:

\begin{equation}
    x = (-1)^s \cdot (1 + m) \cdot 2^{e - 127}
\end{equation}

where:
\begin{equation}
    m = \sum_{i=1}^{23} m_i \cdot 2^{-i}
\end{equation}

\paragraph{Interpretation:}
\begin{itemize}
    \item $(-1)^s$ determines the sign
    \item $(1 + m)$ is the normalized significand
    \item $e - 127$ is the unbiased exponent
\end{itemize}

\paragraph{Special Cases:}
\begin{itemize}
    \item $e = 0$ → denormalized numbers
    \item $e = 255$ → infinity or NaN
\end{itemize}

\subsection{From Raw Memory to Numerical Values}

All numerical values in a program are ultimately stored as sequences of bytes in memory.

A 32-bit float occupies 4 consecutive bytes:

\begin{equation}
    \text{float} \rightarrow [b_0, b_1, b_2, b_3]
\end{equation}

These bytes are interpreted according to the IEEE 754 structure:

\begin{equation}
    \underbrace{s}_{1\ \text{bit}} \quad
    \underbrace{e}_{8\ \text{bits}} \quad
    \underbrace{m}_{23\ \text{bits}}
\end{equation}

\paragraph{Example: Representing $x = 7.98$}

The decimal number $7.98$ cannot be represented exactly in binary, so it is approximated.

Its IEEE 754 single-precision representation is:

\begin{equation}
    x \approx 7.980000019073486
\end{equation}

Binary layout:

\begin{equation}
    0\ 10000001\ 11111111001100110011010
\end{equation}

Breaking this down:

\begin{itemize}
    \item Sign bit:
    \begin{equation}
        s = 0 \quad (\text{positive})
    \end{equation}

    \item Exponent:
    \begin{equation}
        e = 10000001_2 = 129_{10}
    \end{equation}
    \begin{equation}
        e_{\text{unbiased}} = 129 - 127 = 2
    \end{equation}

    \item Mantissa:
    \begin{equation}
        m = 1.1111111001100110011010_2
    \end{equation}
\end{itemize}

Reconstructing the value:

\begin{equation}
    x = (-1)^0 \cdot (1 + m_{\text{fraction}}) \cdot 2^2
\end{equation}


\paragraph{Key Insight:}

There is no intrinsic meaning in memory—only interpretation.

\begin{equation}
    \text{Bits} \xrightarrow{\text{interpretation}} \text{Number}
\end{equation}

The same sequence of bits may represent:
\begin{itemize}
    \item an integer
    \item a floating point number
    \item part of a vertex
    \item or any structured data
\end{itemize}

\paragraph{Conclusion:}

All higher-level constructs (vectors, matrices, geometry) are built from:

\begin{equation}
    \text{Bits} \rightarrow \text{Bytes} \rightarrow \text{Floats} \rightarrow \text{Mathematical Objects}
\end{equation}

\section{Geometry Representation in Memory}

\subsection{Vertices (Position Vectors)}

A \textbf{vertex} represents a point in 3D space.

In its simplest form, a vertex consists of three floating-point values:

\begin{equation}
    \mathbf{v} =
    \begin{bmatrix}
        x \\
        y \\
        z
    \end{bmatrix}
    \in \mathbb{R}^3
\end{equation}

For transformation purposes, vertices are often expressed in homogeneous coordinates:

\begin{equation}
    \mathbf{v}_h =
    \begin{bmatrix}
        x \\
        y \\
        z \\
        1
    \end{bmatrix}
    \in \mathbb{R}^4
\end{equation}

\paragraph{Memory Representation:}

Each coordinate is a 32-bit float:

\begin{equation}
    \text{vertex} = [x, y, z]
\end{equation}

\begin{equation}
    \text{size of vertex} = 3 \times 4\ \text{bytes} = 12\ \text{bytes}
\end{equation}

Thus, a vertex buffer is:

\begin{equation}
    V = [\mathbf{v}_0, \mathbf{v}_1, \dots, \mathbf{v}_n]
\end{equation}


\subsection{Triangles (Faces)}

The fundamental rendering primitive is the \textbf{triangle}.

A triangle is defined by three vertices:

\begin{equation}
    \triangle = (\mathbf{v}_0, \mathbf{v}_1, \mathbf{v}_2)
\end{equation}

Geometrically, any point inside the triangle can be expressed using barycentric coordinates:

\begin{equation}
    \mathbf{p} = \alpha \mathbf{v}_0 + \beta \mathbf{v}_1 + \gamma \mathbf{v}_2
\end{equation}

subject to:

\begin{equation}
    \alpha + \beta + \gamma = 1, \quad \alpha, \beta, \gamma \geq 0
\end{equation}

\paragraph{Memory Representation (Non-Indexed):}

Each triangle stores three full vertices:

\begin{equation}
    \text{triangle size} = 3 \times 12\ \text{bytes} = 36\ \text{bytes}
\end{equation}


\subsection{Indexed vs Non-Indexed Geometry}

\paragraph{Non-Indexed Representation:}

Each triangle explicitly stores its vertices:

\begin{equation}
    T = [\mathbf{v}_0, \mathbf{v}_1, \mathbf{v}_2]
\end{equation}

This leads to duplication of shared vertices.

\paragraph{Indexed Representation:}

Vertices are stored once in a vertex buffer:

\begin{equation}
    V = [\mathbf{v}_0, \mathbf{v}_1, \dots, \mathbf{v}_n]
\end{equation}

Triangles are defined using indices:

\begin{equation}
    F = [(i_0, i_1, i_2), (i_3, i_4, i_5), (i_6, i_7, i_8)]
\end{equation}

Where winding order is used to specify face direction.

Each index refers to a vertex:

\begin{equation}
    \triangle_k = (\mathbf{v}_{i_a}, \mathbf{v}_{i_b}, \mathbf{v}_{i_c})
\end{equation}

\paragraph{Memory Efficiency:}

Let:
\begin{equation}
    n_v = \text{number of vertices}, \quad n_t = \text{number of triangles}
\end{equation}

Then:

\begin{equation}
    \text{non-indexed size} = 36 \cdot n_t
\end{equation}

\begin{equation}
    \text{indexed size} = 12 \cdot n_v + 12 \cdot n_t
\end{equation}

(assuming 32-bit indices)

\subsection{Meshes (Objects)}

A \textbf{mesh} (object) consists of:

\begin{itemize}
    \item A vertex buffer $V$
    \item An index buffer $F$
\end{itemize}

\begin{equation}
    \mathcal{M} = (V, F)
\end{equation}

The mesh is defined in \textbf{local (object) space}, meaning:

\begin{equation}
    \mathbf{v}_{local} \in \mathbb{R}^3
\end{equation}

All vertices are expressed relative to the object's origin:

\begin{equation}
    \mathbf{o}_{local} = (0,0,0)
\end{equation}


\subsection{Scenes as Collections of Objects}

A \textbf{scene} consists of multiple objects:

\begin{equation}
    \mathcal{S} = \{\mathcal{M}_1, \mathcal{M}_2, \dots, \mathcal{M}_k\}
\end{equation}

Each object has an associated transformation:

\begin{equation}
    M_i \in \mathbb{R}^{4 \times 4}
\end{equation}

which maps local coordinates to world coordinates:

\begin{equation}
    \mathbf{v}_{world} = M_i \mathbf{v}_{local}
\end{equation}

\paragraph{Global Interpretation:}

All objects share a common coordinate system:

\begin{equation}
    \mathbf{v}_{world} \in \mathbb{R}^3
\end{equation}

Thus, the scene becomes a unified collection of geometry in world space.

\paragraph{Final Structure in Memory:}

\begin{equation}
    \text{Scene} =
    \big\{
    (V_1, F_1, M_1),
    (V_2, F_2, M_2),
    \dots
    \big\}
\end{equation}

Each component ultimately reduces to:

\begin{equation}
    \text{Bytes} \rightarrow \text{Floats} \rightarrow \text{Vertices} \rightarrow \text{Triangles} \rightarrow \text{Meshes} \rightarrow \text{Scene}
\end{equation}

\section{Coordinate Systems and Spaces}

In computer graphics, geometry is transformed through a sequence of coordinate systems. Each space provides a different frame of reference for interpreting vertex positions.

\begin{equation}
    \mathbf{v}_{local}
    \rightarrow
    \mathbf{v}_{world}
    \rightarrow
    \mathbf{v}_{camera}
    \rightarrow
    \mathbf{v}_{clip}
    \rightarrow
    \mathbf{v}_{ndc}
    \rightarrow
    \mathbf{v}_{screen}
\end{equation}


\subsection{Local (Object) Space}

Local space (or object space) is the coordinate system in which a mesh is originally defined.

\begin{equation}
    \mathbf{v}_{local} =
    \begin{bmatrix}
        x \\
        y \\
        z \\
        1
    \end{bmatrix}
\end{equation}

Key properties:

\begin{itemize}
    \item Origin is the object's center:
    \begin{equation}
        \mathbf{o}_{local} = (0,0,0)
    \end{equation}
    \item Basis vectors define the object's orientation:
    \begin{equation}
        \{\mathbf{e}_x^{local}, \mathbf{e}_y^{local}, \mathbf{e}_z^{local}\}
    \end{equation}
\end{itemize}

Vertices are stored in this space in memory.


\subsection{World Space}

World space is a global coordinate system shared by all objects in the scene.

Each object is placed into world space using a model transformation:

\begin{equation}
    \mathbf{v}_{world} = M_{model} \mathbf{v}_{local}
\end{equation}

where:

\begin{equation}
    M_{model} = T \cdot R \cdot S
\end{equation}

Explicitly:

\begin{equation}
    M_{model} =
    \begin{bmatrix}
        r_{11}S_x & r_{12}S_y & r_{13}S_z & t_x \\
        r_{21}S_x & r_{22}S_y & r_{23}S_z & t_y \\
        r_{31}S_x & r_{32}S_y & r_{33}S_z & t_z \\
        0 & 0 & 0 & 1
    \end{bmatrix}
\end{equation}

Key properties:

\begin{itemize}
    \item Common reference frame for all objects
    \item Fixed global basis:
    \begin{equation}
        \{\mathbf{e}_x^{world}, \mathbf{e}_y^{world}, \mathbf{e}_z^{world}\}
    \end{equation}
\end{itemize}


\subsection{Camera (View) Space}

Camera space (or view space) represents coordinates relative to the camera.

Instead of moving the camera, we transform the world:

\begin{equation}
    \mathbf{v}_{camera} = M_{view} \mathbf{v}_{world}
\end{equation}

The view matrix is the inverse of the camera transform:

\begin{equation}
    M_{view} = (T_{cam} R_{cam})^{-1}
\end{equation}

If the camera has position $\mathbf{c}$ and rotation $R_{cam}$:

\begin{equation}
    M_{view} =
    \begin{bmatrix}
        R_{cam}^T & -R_{cam}^T \mathbf{c} \\
        0 & 1
    \end{bmatrix}
\end{equation}

Key properties:

\begin{itemize}
    \item Camera is located at origin:
    \begin{equation}
        \mathbf{o}_{camera} = (0,0,0)
    \end{equation}
    \item Viewing direction typically along negative z-axis
\end{itemize}


\subsection{Clip Space}

Clip space is obtained by applying the projection matrix:

\begin{equation}
    \mathbf{v}_{clip} = M_{proj} \mathbf{v}_{camera}
\end{equation}

A standard perspective projection matrix is:

\begin{equation}
    M_{proj} =
    \begin{bmatrix}
        \frac{1}{a \tan(\frac{\mathrm{fov}}{2})} & 0 & 0 & 0 \\
        0 & \frac{1}{\tan(\frac{\mathrm{fov}}{2})} & 0 & 0 \\
        0 & 0 & \frac{f+n}{n-f} & \frac{2fn}{n-f} \\
        0 & 0 & -1 & 0
    \end{bmatrix}
\end{equation}

where:

\begin{equation}
    a = \text{aspect ratio}, \quad n = \text{near plane}, \quad f = \text{far plane} \quad \mathrm{fov} = \mathrm{field\ of\ view}
\end{equation}

Vertices in clip space are 4D:

\begin{equation}
    \mathbf{v}_{clip} =
    \begin{bmatrix}
        x_c \\
        y_c \\
        z_c \\
        w_c
    \end{bmatrix}
\end{equation}


\subsection{Normalized Device Coordinates (NDC)}

To obtain normalized device coordinates, perform the perspective divide:

\begin{equation}
    \mathbf{v}_{ndc} =
    \begin{bmatrix}
        x_c / w_c \\
        y_c / w_c \\
        z_c / w_c
    \end{bmatrix}
\end{equation}

This maps coordinates into a canonical cube:

\begin{equation}
    x, y, z \in [-1, 1]
\end{equation}

Key properties:

\begin{itemize}
    \item Removes perspective distortion via division by $w$
    \item Enables uniform clipping and rasterization
\end{itemize}


\subsection{Screen Space}

Screen space maps normalized coordinates to pixel coordinates.

For a viewport of width $W$ and height $H$:

\begin{equation}
    x_{screen} = \frac{W}{2} (x_{ndc} + 1)
\end{equation}

\begin{equation}
    y_{screen} = \frac{H}{2} (1 - y_{ndc})
\end{equation}

\begin{equation}
    z_{screen} = \frac{1}{2}(z_{ndc} + 1)
\end{equation}

Final screen coordinate:

\begin{equation}
    \mathbf{v}_{screen} =
    \begin{bmatrix}
        x_{screen} \\
        y_{screen} \\
        z_{screen}
    \end{bmatrix}
\end{equation}


\subsection{Summary of Transform Pipeline}

The full transformation pipeline is:

\begin{equation}
    \mathbf{v}_{clip} =
    M_{proj} \cdot
    M_{view} \cdot
    M_{model} \cdot
    \mathbf{v}_{local}
\end{equation}

Followed by:

\begin{equation}
    \mathbf{v}_{ndc} = \frac{1}{w} \mathbf{v}_{clip}
\end{equation}

\begin{equation}
    \mathbf{v}_{screen} = \text{Viewport}(\mathbf{v}_{ndc})
\end{equation}

\subsection{Conceptual Interpretation}

Each transformation corresponds to a change of coordinate system:

\begin{equation}
    \text{Local} \rightarrow \text{World} \rightarrow \text{Camera} \rightarrow \text{Clip} \rightarrow \text{Screen}
\end{equation}

At every stage:

\begin{itemize}
    \item Coordinates are expressed relative to a specific basis
    \item Matrix multiplication re-expresses vectors in a new basis
\end{itemize}

Thus, the entire graphics pipeline can be viewed as a sequence of basis transformations applied to vertex data.

\end{document}
